\documentclass[11pt,a4paper]{article}
\usepackage{fontspec,amsmath,amssymb,amsthm,microtype}
\setmainfont{lmroman10-regular.otf}[BoldFont=lmroman10-bold.otf,ItalicFont=lmroman10-italic.otf,BoldItalicFont=lmroman10-bolditalic.otf]
\setmonofont{lmmono10-regular.otf}[Scale=MatchLowercase]
\usepackage[margin=25mm]{geometry}
\usepackage[unicode,hidelinks]{hyperref}
\setlength{\emergencystretch}{3em}
\setlength{\parskip}{0.4em}
\setlength{\parindent}{0pt}
\allowdisplaybreaks[2]
\raggedbottom
\title{\texorpdfstring{The arbitrary-\allowbreak{}shift kinetic contraction for actual L\ensuremath{{}^{2}} random vectors}{The arbitrary-shift kinetic contraction for actual L² random vectors}}
\author{LangevinLean mathematical exposition}
\date{}
\begin{document}\maketitle
\noindent Source module: \texttt{Kinetic\allowbreak{}Langevin/\allowbreak{}Lp\allowbreak{}General\allowbreak{}Contraction.\allowbreak{}lean}. The accompanying HTML pairs this mathematical exposition with the exact Lean statements and records the scope of the formalization.


The continuous mixing theorem permits an arbitrary shifted kinetic metric, not only the convenient zero-shift case used by KLMC. The pointwise strong-smooth gradient sector bound already lifts to L\ensuremath{{}^{2}}. Here it gives the complete shifted energy and position estimates for actual shared-input L\ensuremath{{}^{2}} kinetic solutions.


\section{\texorpdfstring{Prove the full shifted L\ensuremath{{}^{2}} energy decay}{Prove the full shifted L² energy decay}}


Let \(d=\gamma-2a>0\). The exact common-noise integral equations make the difference differentiable, with the original gradient difference lifted pointwise. Lean differentiates the explicit shifted energy and applies the general dissipativity inequality using the genuine lifted sector bound. The resulting rate is \(2\max(a^2-m,M-(\gamma-a)^2)/d\). A right-derivative Grönwall argument integrates this differential inequality, including the starting endpoint. No potential on the L\ensuremath{{}^{2}} space is introduced.



{\raggedright\noindent Lean declarations: \texttt{lp\_\allowbreak{}kinetic\_\allowbreak{}integral\_\allowbreak{}energy\_\allowbreak{}decay\_\allowbreak{}general}.\par}

\section{\texorpdfstring{Project the shifted energy to the position error}{Project the shifted energy to the position error}}


Equal initial velocities make the initial energy equal to \(((\gamma-a)^2+a^2)\|X_0-Y_0\|_2^2\). The invertible transformed coordinates bound the position error by the energy divided by \((\gamma-2a)^2\). Lean combines these identities with the energy decay, proves the exact squared-exponential identity, and takes the nonnegative square-root inequality. The conclusion is the full source-paper prefactor \(\sqrt{2((\gamma-a)^2+a^2)}/(\gamma-2a)\) multiplied by the explicit exponential rate and the actual initial L\ensuremath{{}^{2}} position cost.



{\raggedright\noindent Lean declarations: \texttt{lp\_\allowbreak{}kinetic\_\allowbreak{}integral\_\allowbreak{}position\_\allowbreak{}decay\_\allowbreak{}general}.\par}

\end{document}
